\section{A degree gap sufficient for exponential failure}\label{sec:consequences}

Theorem~A gives a test for an impossibility proof. At a degree cap where
all admissible families have a rate gap $\rho$, a putative transposer
must pay it from its deficit allowance $K f_p(\Delta)$. The machine
constant $K$ is fixed before $\Delta$ is chosen.

\begin{corollary}[A sufficient gap at a sequence of caps]\label{cor:gap}
\formal{TapeStorageRates/Main.lean}{TapeStorageRates.exponentially_small_at_widths_of_sparse_gap}
Fix a tape cap $k\ge0$ and put $p=4k^2+1$. Suppose there are integer
caps $\Delta_j\to\infty$ and real numbers $\rho_j$ such that
$\rho_j/f_p(\Delta_j)\to\infty$. At each $\Delta_j$, assume that
every admissible family has limiting lower rate at most $1-\rho_j$.
For each fixed machine $M$ with at most $k$ work tapes and each $c>0$,
there are $\gamma>0$ and $n_0$ such that
\[
   \frac{|G_{d,b}|}{2^{d^2b}}
      \le 2^{-\gamma d^2b}
      \qquad(d\ge n_0,\ b\ge n_0).
\]
In particular the same exponential upper bound holds at every
sufficiently large canonical dimension, and (T) fails for such machines.
\end{corollary}

\begin{outline}
If the joint near-full success hypothesis of \cref{thm:main} held,
each sufficiently large cap would admit a family of limiting lower rate
at least $1-K f_p(\Delta)$. A machine with fewer than $k$ tapes also
satisfies this bound, since increasing the tape count increases the
deficit allowance at large caps. Choose $j$ with $\Delta_j$ large
enough for the lower bound and $\rho_j>K f_p(\Delta_j)$. The required
family violates the assumed upper bound at that same cap.
Negating the near-full success hypothesis now gives some $\gamma>0$ and
$n_0$ such that, throughout $d,b\ge n_0$, each nonempty success set has
$\log_2|G_{d,b}|\le(1-\gamma)d^2b$. Empty success sets satisfy the
probability bound as well. Finally $b(d)\to\infty$ with $d$.
\end{outline}

For example, suppose $R(\Delta)\le1-\rho(\Delta)$ at every sufficiently
large $\Delta$, with $\rho(\Delta)\ge0$. Writing
$y=\log_2\log_2\Delta$, a gap $\rho(\Delta)\ge c_1/y^a$ for a constant
$c_1>0$ and $0<a<1/p$ suffices for the fixed tape cap $k$.
If $\limsup_{\Delta\to\infty}\rho(\Delta)y^a=\infty$ for every $a>0$,
the exclusion applies to every fixed finite tape count. Indeed
$y^{1/(2p)}f_p(\Delta)\to0$. A gap bounded below by a positive
constant times $1/\log_2 y$ meets this all-tapes condition.
The elementary gap $1/(\Delta+1)$ obtained by ordering coordinates
has ratio to $f_p(\Delta)$ tending to zero
\cite[open problems]{PaperThirteen}. The slowly decreasing gap criterion
for oblivious machines appears after Corollary~O in \cite{PaperFourteen}.
The power comparison and the elementary-gap comparison are checked in
\lean{TapeStorageRates.ratio_tendsto_of_power_gap} and
\lean{TapeStorageRates.elementary_gap_ratio_tendsto_zero}.
The coding upper bound at the scale required for \cref{cor:gap} remains open.
